% ====================================================================
% si_fr.tex — [SI-FR] Si-host 정칙용액(Frumkin, Ω<2RT 고용체) 커널 (신규 subsection)
%   삽입 위치: Chapter 3(Si) §3.2 케이스별 활물질(ch3v22_sec02_cases, \label{sec:si-cases}) 인접.
%   저자 sub W9 (comp_R1). 초안 — 실제 문건 편집 금지, 본 폴더 전용.
%   강조: adversarial completeness + honest limits.
% ====================================================================
\subsection{Si-host 정칙용액 커널 --- 단일상 고용체의 Frumkin peak 과 그 식별 한계}\label{ssec:si-fr}
\S\ref{sec:si-cases} 가 원소 Si$\cdot$SiO$_x$$\cdot$Si--C 세 계열의 개형을 실측으로 읽었으니, 이 소절은 그
개형을 \emph{한} 미분용량 커널로 닫는다. 핵심 주장은 하나다 --- 순환 a-Si(비정질 실리콘)의 리튬화는 흑연 staging 같은
sharp 두-상 전이가 아니라 \emph{단일상 연속 고용체}이며, 그 dQ/dV 는 흑연이 두-상에서 쓰는 델타-펴짐 종이 아니라
\emph{Frumkin(정칙용액) 단일상 peak} 으로 곧장 닫힌다. 흑연 골격(식~\eqref{eq:eqpeak} 평형 peak,
식~\eqref{eq:xieq} logistic)의 이상 극한($\Omega{=}0$)이 로지스틱 종이었듯, $\Omega\ne0$ 의 단일상 일반형이
이 Frumkin 종이다 --- 새 물리가 아니라 같은 정칙용액 자유에너지의 단일상 가지다.

\textbf{(a) 출발 --- 실측 판정: a-Si 는 단일상 고용체.} 두-상인지 고용체인지는 등온선 폭이 가른다. 우리 정칙용액
피팅 진단(regsol\_si)이 a-Si 전이들의 폭을 열 척도로 정규화한 값이
\begin{equation}
\frac{w_j^{\mathrm{Si}}}{RT/F}\;=\;[\,1.45,\;2.74,\;1.09\,]\;\gtrsim\;1,
\label{eq:sifr-width}
\end{equation}
곧 자연 열폭 $RT/F$ 의 $\sim$1--3 배로 \emph{넓다}. 이는 흑연 두-상 전이의 평형 등온선 폭(near-delta plateau,
$w/(RT/F)\!\ll\!0.12$)과 $\sim$20--50$\times$ 대조된다 --- 두-상이면 평형이 예측하는 것이 폭이 아니라 날카로운
선(\S\ref{sec:width} 이중지위 둘째 지위)인데, a-Si 는 정반대로 \emph{평형이 유한 폭을 예측}한다. 이 폭 대조가
``a-Si $=$ 단일상 고용체''의 실측 근거이며, 미시적으로는 비정질 host 의 자리 에너지 분산이 연속 고용체를 만드는
것과 정합한다\cite{chevrier_dahn2009,artrith2018}.

\textbf{(b) 연산 --- Frumkin 단일상 peak 의 유도.} 단일상 고용체의 전위-조성 관계는 정칙용액 화학퍼텐셜을 전위로
환산해 얻는다. 자리 점유 $\theta$(리튬화 진행)에 대해
\begin{equation}
V(\theta)=U^\circ-\frac{RT}{F}\ln\frac{\theta}{1-\theta}-\frac{\Omega}{F}\,(1-2\theta)
\label{eq:sifr-V}
\end{equation}
이고($\theta{=}1{-}\xi$ 여집합 좌표라 흑연 식~\eqref{eq:lco-Veq} 의 $\xi$-형과 부호까지 동치 --- 로그 몫과
$(1-2\theta)$ 몫이 여집합 교환에서 각각 부호를 바꿔 상쇄), 이를 $\theta$ 로 미분하면
\begin{equation}
\frac{\dd V}{\dd\theta}=-\frac{RT}{F}\cdot\frac{1}{\theta(1-\theta)}+\frac{2\Omega}{F}
=-\frac{1}{F}\Big[\frac{RT}{\theta(1-\theta)}-2\Omega\Big].
\label{eq:sifr-dVdtheta}
\end{equation}
전이 용량 $Q$ 에 대해 $\dd Q=Q\,\dd\theta$ 이므로 미분용량은 연쇄율 $\dd Q/\dd V=(\dd Q/\dd\theta)/(\dd V/\dd\theta)$
로 닫히고, 절댓값을 취하면(peak 은 방향 불변 양수)
\begin{equation}
\boxed{\;\frac{\dd Q}{\dd V}=\frac{Q\,F}{\big|\,RT/[\theta(1-\theta)]-2\Omega\,\big|},
\qquad \Omega<2RT\ \ (\text{단일상 고용체})\;.}
\label{eq:sifr-kernel}
\end{equation}
이것이 Si-host 전이 하나의 평형 종이다 --- 흑연 평형 peak 식~\eqref{eq:eqpeak}($Q_j\xi(1-\xi)/w_j$)의 이상 극한과
같은 대상의 $\Omega\ne0$ 일반형이며, $\Omega{\to}0$ 에서 분모 $RT/[\theta(1-\theta)]$ 가 되어 로지스틱 종
$QF\theta(1-\theta)/RT$ 로 정확히 환원된다(미지정 커널의 로지스틱 폴백과 bit-exact --- 시드표 코드 규약).

\begin{signbox}
\textbf{$\Omega$ 의 세 지위(부호$\cdot$극한).} 분모 $RT/[\theta(1-\theta)]-2\Omega$ 는 $\theta{=}\tfrac12$ 에서
최소 $4RT-2\Omega$ 다. (i) $0<\Omega<2RT$: 분모가 전 구간 양수라 단일 유한 peak(단일상). $\Omega$ 가 $2RT$ 로
다가갈수록 peak 정점이 $Q F/(4RT-2\Omega)$ 로 \emph{솟고}(sharpen), $\Omega{\to}2RT^-$ 에서 발산 --- 이것이 두-상
문턱이다(그 위는 miscibility gap, \S\ref{ssec:si-fr} 밖). (ii) $\Omega{=}0$: 이상 로지스틱 종($RT/F$ 폭). (iii)
$\Omega<0$(질서화 경향): $-2\Omega>0$ 이 분모를 키워 peak 가 \emph{낮고 넓어짐}(broaden). 곧 하나의 $\Omega$ 가
peak 의 날카로움을 연속으로 조율하며, $\Omega<2RT$ 라는 부등호가 ``새 상 없이 한 종''임을 보장한다.
\end{signbox}

\textbf{(c) 중간식 --- Si-host 는 소수의 넓은 갤러리, 유일 두-상은 $c$-Li$_{15}$Si$_4$.} 연속 고용체는 하나의
매끄러운 $U(\theta)$ 곡선이지만 forward 모델은 이를 \emph{소수의 넓은 Frumkin 종}(식~\eqref{eq:sifr-kernel})의
합으로 이산화해 근사한다 --- 각 종이 넓어($\Omega<2RT$) 겹치며 연속 경사를 재구성한다(케이스별 개형은
표~\ref{tab:si-cases}$\cdot$그림~\ref{fig:si-cases-shape}). 이 그림에서 \emph{예외}는 하나뿐이다: 결정질 Si 의
\emph{1차} 심리튬화 끝 $\sim$50--60 mV 의 $c$-Li$_{15}$Si$_4$ 결정화 plateau 는 진짜 두-상 전이라\cite{obrovac_christensen2004,ogata_nmr2014}
Frumkin 단일상 커널의 대상이 아니며(별도 sharp feature), 순환 a-Si 경로에는 \emph{부재}한다\cite{chevrier_dahn2009}.
곧 순환 Si-host 는 ``넓은 고용체 갤러리 $+$ (1차에 한한) 단일 두-상 결정화''로 갈리며, 본 커널은 앞쪽을 담고
뒤쪽은 first-cycle 전용 feature 로 분리한다(honest scoping).

\textbf{(d) 박스 --- 블렌드 가산 중첩의 정당성.} Si-host 종을 흑연 종과 같은 전위축에 얹는 것(블렌드
식~\eqref{eq:blend-dqdv})은 두 host 가 \emph{같은} 전위 손잡이(공통 $\mu$, \S\ref{sec:si-blend})를 공유하기
때문이며, 평형에서 관측 미분용량은 두 host 기여의 배경 위 선형 합이다:
\begin{equation}
\boxed{\;\frac{\dd Q}{\dd V}\bigg|_{U_\mathrm{oc}}=C_\bg+\sum_{j}\underbrace{\frac{Q_j^{\mathrm{gr}}\,\xi_{\eq,j}^{\mathrm{gr}}(1-\xi_{\eq,j}^{\mathrm{gr}})}{w_j^{\mathrm{gr}}}}_{\text{흑연 종}}
+\sum_{k}\underbrace{\frac{Q_k^{\mathrm{Si}}\,F}{\big|RT/[\theta_k(1-\theta_k)]-2\Omega_k^{\mathrm{Si}}\big|}}_{\text{Si-host Frumkin 종(식~\eqref{eq:sifr-kernel})}}\;}
\label{eq:sifr-blend}
\end{equation}
--- 두 host 의 서로 다른 커널(로지스틱 종 vs Frumkin 종)이 같은 $U_\mathrm{oc}$ 축을 공유해 합산된다. 이 가산
중첩이 유효함은 환원차수 블렌드 모형이 혼합 dQ/dV 를 두 성분의 ``분명한 중첩(clearly a superposition)''으로
읽은 것과 정합한다\cite{tu_blend2024}(공통-$\mu$ 실증 계보 --- Verbrugge MSMR\cite{verbrugge_lisi2016}). 다만 이
중첩은 \emph{평형} 유효이고 유한 율속 블렌드의 비가산 편차는 GS-2 공백(\S\ref{ssec:si-blend-gs2})이 별도로 진단한다.

\medskip
\textbf{★적대적 검토 대응 --- 세 물음.} \emph{(1) ``a-Si 의 넓은 peak 은 두-상을 못 맞춘 변명 아닌가.''}
아니다 --- 폭 판정~\eqref{eq:sifr-width} 이 정반대 방향의 실측이다: 두-상이면 평형 폭이 $\ll RT/F$(델타)여야
하는데 a-Si 는 $\gtrsim RT/F$(20--50$\times$ 초과)라, 넓은 종은 커널의 실패가 아니라 \emph{단일상이라는 실측의
직접 결과}다. 또한 제일원리$\cdot$기계학습 퍼텐셜 계산이 a-Li$_x$Si 의 연속 경사 $U(x)$ 를 독립 재현한다\cite{chevrier_dahn2009,artrith2018}.
\emph{(2) ``$\Omega_\mathrm{Si}$ 는 식별 가능한가.''} \textbf{제한적이다} --- 물리가 고정하는 것은 \emph{영역}
($\Omega<2RT$, 단일상)뿐이고 \emph{점값}이 아니다. 아래 정직한 한계 참조. \emph{(3) ``Frumkin 커널이 없어도
로지스틱으로 되지 않나.''} $\Omega{=}0$ 로지스틱은 넓은 대칭 종만 내지만, a-Si 종은 $\Omega\ne0$ 로 sharpen
/broaden 하는 \emph{비대칭 조율}이 필요하다(케이스별 개형 차, 그림~\ref{fig:si-cases-shape}) --- 커널 분기가 이
자유도 하나를 열되, 미지정 시 로지스틱 폴백으로 bit-exact 회수돼 흑연 골격을 손상하지 않는다.

\begin{warnbox}
\textbf{정직한 한계.} (i) \emph{$\Omega_\mathrm{Si}$ 는 범위 시드일 뿐 점-식별 안 됨.} 시드는 $\Omega<2RT$
(넓은 고용체 지향, $\delta$-물리캡 바닥 $\sim0.2RT$)라는 \emph{범위}이고, 실제 점값은 피팅에 위임한다 --- 넓은
단일상 종은 $\Omega$ 에 대한 민감도가 낮아(peak 정점이 $\Omega{\to}2RT$ 근처에서만 급변) 상온 단일 곡선만으로는
$\Omega_\mathrm{Si}$ 를 잘 조이지 못한다. 물리가 확정하는 것은 \emph{상성격}(단일상)이지 $\Omega$ 의 숫자가
아니다. (ii) \emph{폭 진단값의 tier.} 식~\eqref{eq:sifr-width} 의 $[1.45,2.74,1.09]$ 는 우리 regsol\_si 피팅
산출이지 문헌 anchor 가 아니다(전이별 문헌 $\Omega_\mathrm{Si}$ 미발견). (iii) \emph{케이스 개형은 시연값.}
표~\ref{tab:si-cases}$\cdot$그림~\ref{fig:si-cases-shape} 의 Si 전이 중심$\cdot$폭은 tier-C 시연값(피팅 override
전제)이며, 특히 SiO$_x$ 절대 전위$\cdot$히스 mV 는 순수 SiO$_x$ 1차 문헌 미특정으로 ``확인 필요'' 공백이다.
(iv) \emph{$\omega$-형식 대응.} 식~\eqref{eq:sifr-kernel} 은 substitutional 열역학 모델의 $\omega$(상호작용) 형식
\cite{verbrugge2017} 과 같은 정칙용액 골격이나, 그 원 논문의 다-사이트 파라미터화와 본 커널의 전이별
$\Omega_k^\mathrm{Si}$ 매핑은 피팅 단계에서 확정한다(구조 대응이지 값 1:1 아님).
\end{warnbox}

\begin{keybox}
\textbf{Si-host 커널 다리.} a-Si 는 폭 판정~\eqref{eq:sifr-width}($w/(RT/F)\gtrsim1$, 흑연 두-상 대비 20--50$\times$)로
\emph{단일상 고용체}이며, 그 평형 peak 은 Frumkin 종~\eqref{eq:sifr-kernel}
$\dd Q/\dd V=QF/|RT/[\theta(1-\theta)]-2\Omega|$($\Omega<2RT$)로 닫힌다 --- $\Omega{\to}2RT$ sharpen$\cdot\Omega{=}0$
로지스틱 폴백(bit-exact)$\cdot\Omega<0$ broaden. 유일 두-상은 1차 $c$-Li$_{15}$Si$_4$ plateau(순환 a-Si 엔 부재).
블렌드 가산 중첩~\eqref{eq:sifr-blend} 은 공통-$\mu$ 로 유효(평형 극한). \emph{정직 한계}: $\Omega_\mathrm{Si}$ 는
범위 시드($\Omega<2RT$)일 뿐 점-식별되지 않고 피팅에 위임된다.
\end{keybox}

% 자산(comp_R1 W9 초안): [W9-SIFR-1 폭 판정 w/(RT/F)≳1 단일상 실측 eq:sifr-width(흑연 두-상 20-50× 대조)]
%   [W9-SIFR-2 Frumkin 단일상 peak 유도 eq:sifr-V→eq:sifr-dVdtheta→boxed eq:sifr-kernel(Ω→0 로지스틱 폴백)]
%   [W9-SIFR-3 Ω 세 지위 signbox(sharpen/폴백/broaden)] [W9-SIFR-4 소수 넓은 갤러리·유일 두-상 c-Li15Si4]
%   [W9-SIFR-5 블렌드 가산중첩 boxed eq:sifr-blend(공통-μ·tu_blend2024)] [W9-SIFR-6 적대 3문+정직 한계 warnbox(Ω 범위시드·점식별 불가)]
